\documentclass[../../../include/open-logic-section]{subfiles}

\begin{document}

\olfileid{sth}{ordinals}{basic}
\olsection{Sipat Dhasar Ordinal}

Kita wis ndeleng yen sawetara ordinal kapisan iku
wilangan asli. Alasan utama ngembangake teori ordinal
yaiku nggedhekake prinsip induksi kang lumaku ing
wilangan asli. Kita bakal tekan kana liwat sawetara
asil dhasar.

\begin{lem}\ollabel{ordmemberord}
Saben !!{element} saka ordinal iku uga ordinal.
\end{lem}

\begin{proof}
Upamane $\alpha$ ordinal kanthi $b \in \alpha$.
Amarga $\alpha$ transitif, $b \subseteq \alpha$.
Mula $\in$ ngurutake $b$ kanthi becik kaya dene
$\in$ ngurutake $\alpha$ kanthi becik.

Kanggo nuduhake yen $b$ transitif, upamane
$x \in c \in b$. Mula $c \in \alpha$ amarga
$b \subseteq \alpha$. Maneh, amarga $\alpha$
transitif, $c \subseteq \alpha$, mula $x \in \alpha$.
Dadi $x, c, b \in \alpha$. Nanging $\in$
ngurutake $\alpha$ kanthi becik, mula $\in$
minangka relasi transitif ing $\alpha$ miturut
\olref[sth][ordinals][wo]{wo:strictorder}. Amarga
$x \in c \in b$, kita entuk $x \in b$. Kanthi
umum, $c \subseteq b$.
\end{proof}

\begin{cor}\ollabel{ordissetofsmallerord}
$\alpha = \Setabs{\beta \in \alpha}{\beta \text{ iku ordinal}}$,
kanggo sembarang ordinal~$\alpha$.
\end{cor}

\begin{proof}
Langsung ndherek saka \olref{ordmemberord}.
\end{proof}

Inti saka rong asil utama sabanjure,
\olref{ordinductionschema} lan \olref{ordtrichotomy},
yaiku ordinal dhewe diurutake kanthi becik dening
relasi keanggotaan:

\begin{thm}[Induksi Transfinit]\ollabel{ordinductionschema}
Kanggo sembarang rumus $\phi(x)$:
\[
	\text{yen }\exists \alpha \phi(\alpha)\text{, mula }\exists \alpha(\phi(\alpha)
	\land  (\forall \beta \in \alpha) \lnot \phi(\beta))
\]
ing kene kuantor-kuantor kang ditampilake kanthi
ora langsung diwatesi marang ordinal.
\end{thm}

\begin{proof}
%\olref{ordinductionschema1}.
Upamane $\phi(\alpha)$ kanggo sawijining ordinal
$\alpha$. Yen $ (\forall \beta \in \alpha)
\lnot \phi(\beta)$, bukti rampung. Yen ora, amarga
$\alpha$ ordinal, anggota-anggotane kang nyukupi
predikat iki mbentuk himpunan bagean ora kosong
kanthi Pamisahan. Himpunan bagean iki nduweni
!!{element} paling cilik miturut $\in$ kang nyukupi
$\phi$, lan unsur iki ordinal miturut
\olref{ordmemberord}. Amarga ordinal wiwitan transitif,
saben ordinal sadurunge unsur kapilih iki uga dadi
anggotane; mula ora ana kang nyukupi predikat kasebut.
%	\olref{ordinductionschema2}. Upamane ana ordinal
%	$\gamma$ kanthi $\lnot\phi(\gamma)$. Mula miturut
%	\olref{ordinductionschema1}, ana ordinal minimal miturut $\in$
%	$\alpha$ kang nyukupi $\lnot\phi(\alpha)$. Dadi $(\forall \beta <
%	\alpha) \phi(\beta)$, saengga anteseden kondisional iku
%	luput.
\end{proof}
\noindent
Elinga yen \olref{ordinductionschema} uga bisa
dinyatakake minangka skema iki:
\[
\text{yen }\forall \alpha((\forall \beta \in \alpha)\phi(\beta) \lif 
\phi(\alpha))\text{, mula }\forall \alpha\phi(\alpha)
\]
yaiku kanthi njupuk $\lnot\phi(\alpha)$ ing
\olref{ordinductionschema}, banjur nindakake
owah-owahan logis kang prasaja.

\begin{thm}[Trikotomi]\ollabel{ordtrichotomy}
$\alpha \in \beta \lor \alpha = \beta \lor \beta \in \alpha$,
kanggo sembarang ordinal $\alpha$ lan $\beta$.
\end{thm}

\begin{proof}
Bukti iki nganggo induksi kaping pindho, yaiku
nggunakake \olref{ordinductionschema} kaping pindho.
Kita ngarani $x$ \emph{bisa dibandhingake} karo
$y$ yen lan mung yen $x \in y \lor x = y \lor y \in x$.

Kanggo induksi, upamane saben ordinal ing~$\alpha$
bisa dibandhingake karo \emph{saben} ordinal. Kanggo
induksi kapindho, upamane $\alpha$ bisa dibandhingake
karo saben ordinal ing~$\beta$. Kita bakal nuduhake
yen $\alpha$ bisa dibandhingake karo~$\beta$.
Miturut induksi ing~$\beta$, mula $\alpha$ bisa
dibandhingake karo saben ordinal; banjur miturut
induksi ing~$\alpha$, \emph{saben} ordinal bisa
dibandhingake karo \emph{saben} ordinal, kaya kang
dikarepake. Cukup upamane $\alpha \notin \beta$
lan $\beta \notin \alpha$, banjur nuduhake yen
$\alpha = \beta$.

Kanggo nuduhake $\alpha \subseteq \beta$, pilihen
$\gamma \in \alpha$; iki ordinal miturut
\olref{ordmemberord}. Miturut pangira-ira induksi
kapisan, $\gamma$ bisa dibandhingake karo $\beta$.
Nanging yen $\gamma = \beta$ utawa
$\beta \in \gamma$, mula $\beta \in \alpha$
(yen perlu, nganggo transitivitas $\alpha$),
nentang pangira-ira kita; mula $\gamma \in \beta$.
Kanthi umum, $\alpha \subseteq \beta$.

Penalaran kang padha persis, nganggo pangira-ira
induksi kapindho, nuduhake $\beta \subseteq \alpha$.
Mula $\alpha = \beta$.
\end{proof}\noindent Mula kadhang kita nulis
$\alpha <\beta$ tinimbang $\alpha \in \beta$,
amarga $\in$ tumindak minangka relasi urutan.
Ora ana alasan kang luwih jero kajaba wis lumrah,
lan amarga luwih gampang nulis $\alpha \leq \beta$
tinimbang $\alpha \in \beta \lor \alpha = \beta$.
\footnote{Kita bisa nulis
$\alpha \mathrel{\underline{\in}} \beta$,
nanging cara iki babar pisan ora lumrah.}

Iki rong akibat langsung saka asil-asil mau;
sing kapisan nggunakake notasi anyar kita:

\begin{cor}\ollabel{ordordered}
Yen $\exists \alpha\phi(\alpha)$, mula
$\exists \alpha(\phi(\alpha)
\land \forall \beta(\phi(\beta) \lif \alpha \leq \beta))$.
Kajaba iku, kanggo sembarang ordinal $\alpha, \beta, \gamma$,
lumaku $\alpha \notin \alpha$ lan
$\alpha \in \beta \in \gamma \lif \alpha \in \gamma$.
\end{cor}

\begin{proof}
Kaya ing \olref[wo]{wo:strictorder}.
\end{proof}

\begin{prob}
Rampungna ``penalaran kang padha persis'' ing bukti
\olref[sth][ordinals][basic]{ordtrichotomy}.
\end{prob}

\begin{cor}\ollabel{corordtransitiveord}
$A$ iku ordinal yen lan mung yen $A$ iku himpunan
transitif kang anggotane ordinal.
\end{cor}

\begin{proof}
\emph{Kiwa menyang tengen.} Miturut
\olref{ordmemberord}. \emph{Tengen menyang kiwa.}
Yen $A$ himpunan transitif kang anggotane ordinal,
mula $\in$ ngurutake $A$ kanthi becik miturut
\olref{ordinductionschema} lan \olref{ordtrichotomy}.
\end{proof}

Saiki, kita wis nerangake
\olref{ordinductionschema} lan \olref{ordtrichotomy}
kaya-kaya loro asil iku negesake yen $\in$ ngurutake
ordinal kanthi becik. Nanging kita kudu
\emph{ngati-ati banget} marang pratelan kaya mangkene,
amarga asil sabanjure:

\begin{thm}[Paradoks Burali-Forti]\ollabel{buraliforti}
Ora ana himpunan kang ngemot kabeh ordinal.
\end{thm}

\begin{proof}
Kanggo bukti kanthi kontradiksi, upamane $O$ iku
himpunan kabeh ordinal. Yen $\alpha \in \beta \in O$,
mula $\alpha$ ordinal miturut \olref{ordmemberord},
dadi $\alpha \in O$. Mula $O$ transitif, saengga
$O$ ordinal miturut \olref{corordtransitiveord}.
Akibate $O \in O$, nentang \olref{ordordered}.
\end{proof}
\noindent
Asil iki dijenengi miturut \citeauthor{Burali-Forti1897}.
Nanging Cantor ing taun 1899---liwat layang marang
Dedekind---kang sepisanan ndeleng kanthi cetha
\emph{kontradiksi} saka pangira-ira yen ana himpunan
kabeh ordinal. Kaya kang diterangake van Heijenoort:
\begin{quote}
  Burali-Forti dhewe nganggep kontradiksi iku mbuktekake,
  liwat \emph{reductio ad absurdum}, yen urutan alami
  ordinal iku mung urutan parsial.
  \citep[p.~105]{Heijenoort1967}
\end{quote}
Yen kalirune Burali-Forti disisihake, asil-asil mau
bisa diringkes mangkene. Saben ordinal iku himpunan
kang diurutake kanthi becik dening keanggotaan;
ordinal-ordinal uga diurutake kanthi becik dening
keanggotaan minangka sakabehan, sanadyan sakabehan
iku ora mbentuk siji himpunan.

Kanggo ngrampungake bab iki, iki sawetara sipat
dhasar ordinal liyane kang ndherek saka
\olref{ordinductionschema} lan \olref{ordtrichotomy}.

\begin{prop}
Saben urutan ordinal kang mudhun kanthi ketat mung
nduweni anggota kang cacahé winates.
\end{prop}

\begin{proof}
Sembarang urutan ordinal tanpa wates kang mudhun
kanthi ketat $\alpha_0 > \alpha_1 > \alpha_2 > \ldots$
ora nduweni anggota minimal miturut $<$,
nentang \olref{ordinductionschema}.
\end{proof}

\begin{prop}\ollabel{ordinalsaresubsets}
$\alpha \subseteq \beta \lor \beta \subseteq \alpha$,
kanggo sembarang ordinal $\alpha, \beta$.
\end{prop}

\begin{proof}
Yen $\alpha \in \beta$, mula $\alpha \subseteq \beta$
amarga $\beta$ transitif. Semono uga, yen
$\beta \in \alpha$, mula $\beta \subseteq \alpha$.
Yen $\alpha = \beta$, mula $\alpha \subseteq \beta$
lan $\beta \subseteq \alpha$. Dadi bukti rampung
miturut \olref{ordtrichotomy}.
\end{proof}

\begin{prop}\ollabel{ordisoidentity}
$\alpha = \beta$ yen lan mung yen
$\ordeq{\alpha}{\beta}$, kanggo sembarang ordinal
$\alpha, \beta$.
\end{prop}

\begin{proof}
Ordinal iku urutan becik; mula iki langsung ndherek
saka Trikotomi (\olref[basic]{ordtrichotomy}) lan
\olref[iso]{wellordnotinitial}.
\end{proof}

\begin{prob}\ollabel{probunionordinalsordinal}
Buktikna yen saben anggota $X$ iku ordinal,
mula $\bigcup X$ uga ordinal.
\end{prob}

\end{document}
